% !TEX program = xelatex
\documentclass[9pt, a4paper]{article}

% =========================================================
% 1. COMPATIBILIDAD Y CODIFICACIÓN
% =========================================================
\usepackage{fontspec}
\setmainfont{DM Sans}        % cuerpo (DM Sans)
\setmonofont{IBM Plex Mono}[  % títulos en monoespaciado (vía \ttfamily)
    Path=/home/sanaruca/.fonts/,
    UprightFont=IBMPlexMono-Regular.ttf,
    BoldFont=IBMPlexMono-SemiBold.ttf,
    ItalicFont=IBMPlexMono-Italic.ttf,
    BoldItalicFont=IBMPlexMono-SemiBoldItalic.ttf,
]  % el "bold" del mono es SemiBold (más fino que Bold)
\newfontfamily\PlexSans{IBM Plex Sans}[
    Path=/home/sanaruca/.fonts/,
    UprightFont=IBMPlexSans-Regular.ttf,
    BoldFont=IBMPlexSans-Bold.ttf,
    ItalicFont=IBMPlexSans-Italic.ttf,
    BoldItalicFont=IBMPlexSans-BoldItalic.ttf,
]  % columna izquierda de Experiencia

% Pesos adicionales de IBM Plex Sans para la columna izquierda de Experiencia
\newfontfamily\PlexMedium{IBM Plex Sans}[
    Path=/home/sanaruca/.fonts/,
    UprightFont=IBMPlexSans-Medium.ttf,
]
\newfontfamily\PlexLight{IBM Plex Sans}[
    Path=/home/sanaruca/.fonts/,
    UprightFont=IBMPlexSans-Light.ttf,
    ItalicFont=IBMPlexSans-LightItalic.ttf,
]
\newfontfamily\PlexExtraLight{IBM Plex Sans}[
    Path=/home/sanaruca/.fonts/,
    UprightFont=IBMPlexSans-ExtraLight.ttf,
]
\newfontfamily\PlexMonoMedium{IBM Plex Mono}[  % medium (500) para la línea T.S.U.
    Path=/home/sanaruca/.fonts/,
    UprightFont=IBMPlexMono-Medium.ttf,
    ItalicFont=IBMPlexMono-MediumItalic.ttf,
]
% DM Sans para la parte derecha del encabezado (texto Light, viñetas Medium)
\newfontfamily\DMSansLight{DM Sans}[
    Path=/home/sanaruca/.fonts/,
    UprightFont=DMSans-Variable.ttf,
    UprightFeatures={RawFeature={wght=300}},
]
\newfontfamily\DMSansMedium{DM Sans}[
    Path=/home/sanaruca/.fonts/,
    UprightFont=DMSans-Medium.ttf,
]
% babel[spanish] requiere texlive-lang-spanish (no instalado). Se omite el idioma.
% Con XeLaTeX los acentos son UTF-8 nativo; no se usa inputenc/fontenc.
% \usepackage[spanish]{babel}

% =========================================================
% 2. CONFIGURACIÓN DE PÁGINA (Layout Dinámico)
% =========================================================
\usepackage{geometry}
\geometry{left=0.9cm, right=0.9cm, top=1.3cm, bottom=0.5cm}

\usepackage{xcolor}
\usepackage{colortbl}
\usepackage{enumitem}
\usepackage{multicol}
\usepackage{fontawesome5}
\usepackage{hyperref}

% PAQUETE CLAVE PARA TU LAYOUT [SECCION_TITULO] | [INFO]
\usepackage{paracol} 

% =========================================================
% CONFIGURACIÓN DE COLORES Y ENLACES
% =========================================================
\definecolor{CVPrimary}{RGB}{0, 0, 0}   % Azul oscuro
\definecolor{CVText}{RGB}{50, 50, 50}       % Texto general
\definecolor{CVLine}{RGB}{217, 217, 217}    % Gris para las líneas separadoras (|)
\definecolor{CVBullet}{RGB}{179, 179, 179}  % Gris para las viñetas (guiones)

\hypersetup{
    colorlinks=true,
    linkcolor=CVPrimary,
    urlcolor=CVPrimary,
    pdftitle={Samir N. Ruza - Desarrollador Full-Stack},
    pdfauthor={Samir N. Ruza},
    pdfsubject={Curriculum Vitae - Samir N. Ruza},
    pdfkeywords={Curriculum Vitae, Samir Ruza, Desarrollador Full-Stack, Full Stack Developer, Node.js, TypeScript, Dart, Flutter, Golang, PHP, SQL, JavaScript, HTML, CSS, Software Development, Desarrollo de Software}
}

\setlength{\parindent}{0pt} % Eliminar sangrías

% =========================================================
% 3. COMANDOS PERSONALIZADOS DEL NUEVO LAYOUT
% =========================================================

% Comando para los Títulos de Sección en la columna izquierda
\newcommand{\CVSection}[1]{
    \vspace{0.1cm}
    \noindent{\large\bfseries\ttfamily\color{CVPrimary} \raggedright \MakeUppercase{#1}\par}
    \vspace{0.5cm}
}

% Comando para los Títulos de Sección en el modo "Normal" (ancho completo)
\newcommand{\CVSectionNormal}[1]{
    \vspace{1cm}
    \noindent{\large\bfseries\ttfamily\color{CVPrimary} \MakeUppercase{#1}}\par
    \vspace{0.5cm}
}

% --- NUEVOS COMANDOS CON LÍNEA SEGMENTADA ---

% Comando para texto simple (como el perfil) con línea lateral
\newcommand{\CVTextBlock}[1]{
    \noindent
    \textcolor{CVLine}{\vrule width 1pt} \hspace{1cm}
    \begin{minipage}[t]{0.95\linewidth}
        \raggedright {\DMSansLight #1}
    \end{minipage}
    \vspace{0.3cm}
}

% Comando para la Experiencia Laboral con línea lateral
\newcommand{\CVEntry}[4]{
    \noindent
    \textcolor{CVLine}{\vrule width 1pt} \hspace{1cm}
    \begin{minipage}[t]{0.95\linewidth}
        \textbf{#1} \hfill \textbf{\color{CVPrimary}#3} \\
        \textit{#2} \vspace{0.1cm} \\
        \textcolor{CVText}{#4}
    \end{minipage}
    \vspace{0.4cm}
}

% Comando para Educación con línea lateral
\newcommand{\CVEdu}[3]{
    \noindent
    \textcolor{CVLine}{\vrule width 1pt} \hspace{1cm}
    \begin{minipage}[t]{0.95\linewidth}
        \textbf{#1} \hfill \textbf{\color{CVPrimary}#3} \\
        \textit{#2}
    \end{minipage}
    \vspace{0.3cm}
}

% Comando para Referencias
\newcommand{\CVRef}[4]{
    \noindent
    \textcolor{CVLine}{\vrule width 1pt} \hspace{1cm}
    \begin{minipage}[t]{0.42\textwidth}
        {\PlexMedium #1} \\
        {\PlexSans #2} \\
        {\PlexLight #3} \\
        {\DMSansMedium tlf.}{\DMSansLight\ #4}
    \end{minipage}
}

% --- Bloque de metadatos en la columna izquierda de Experiencia ---
% Orden exacto: cargo / empresa / lugar / periodo (uno bajo otro)
\newcommand{\CVExpMeta}[4]{
    {\noindent
    {\PlexMedium\color{CVPrimary} #1} \\
    {\PlexSans #2} \\
    {\PlexLight\itshape #3} \\
    {\PlexExtraLight #4}}
    \vspace{0.85cm}
}

% --- Cuerpo de descripción en la columna derecha de Experiencia ---
\newcommand{\CVExpBody}[1]{
    \noindent
    \textcolor{CVLine}{\vrule width 1pt} \hspace{1cm}
    \begin{minipage}[t]{0.92\linewidth}
        \raggedright {\DMSansLight #1}
    \end{minipage}
    \vspace{0.8cm}
}

% --- Dato de contacto: abreviatura de ancho fijo (Medium) + valor (Light) ---
\newcommand{\CVContact}[2]{
    {\DMSansMedium\makebox[1.1cm][l]{#1}} {\DMSansLight #2} \\[0.04cm]
}

% --- Ítem de Educación: título / ente emisor / lugar / fecha (uno bajo otro) ---
\newcommand{\CVEdItem}[4]{
    \noindent
    {\PlexMedium #1} \\
    {\PlexSans #2} \\
    {\PlexLight\itshape #3} \\
    {\PlexExtraLight #4}
    \vspace{0.85cm}
}


% =========================================================
% ENCABEZADO REUTILIZABLE (dos columnas con línea gris)
% La regla '|' llega hasta el borde superior de la hoja mientras
% que el nombre/contacto quedan desplazados hacia abajo.
% =========================================================
\newcommand{\CVHeader}{
    {\topskip=0pt
    \vspace*{-1.3cm}
    \arrayrulecolor{CVLine}
    \noindent
    \begin{tabular}{@{}p{\dimexpr0.74\linewidth-0.6cm\relax}@{\hspace{0.4cm}}|@{\hspace{0.8cm}}p{\dimexpr0.26\linewidth-0.2cm\relax}@{}}
    \rule{0pt}{2.3cm}
    {\Huge\ttfamily\hspace*{-2.9pt}\addfontfeature{LetterSpace=30}\textbf{\MakeUppercase{Samir N. Ruza}}} \\[-5pt]
    {\large\color{CVPrimary}\PlexMonoMedium T.S.U. en Informática para la Gestión Social} \\[-5pt]
    {\large\color{CVText}\ttfamily Desarrollador Full-Stack}
    &
    \footnotesize
    \begin{minipage}[t]{\linewidth}
    \vspace*{-1.8cm}
    \raggedright
    \CVContact{tlf.}{+58 424 9088 629}
    \CVContact{ce.}{sanaruca@gmail.com}
    \CVContact{github}{/Sanaruca}
    \end{minipage}
    \end{tabular}}
}

\begin{document}

% =========================================================
% ENCABEZADO: [NOMBRE] / [TITULO] / [Cargo - Contacto]
% =========================================================
% ENCABEZADO (página 1)
\CVHeader
\vspace{1cm}

% =========================================================
% INICIO DEL LAYOUT DIVIDIDO
% =========================================================
\columnratio{0.22}               % 22% de ancho para la izquierda, 78% para la derecha
\setlength{\columnsep}{1.6cm}    % Espacio entre el texto y la columna izquierda
% AQUÍ ESTÁ EL CAMBIO PRINCIPAL: Se borraron las líneas que hacían la raya global

\begin{paracol}{2}

% --- PERFIL PROFESIONAL ---
\CVSection{Perfil Profesional}
\switchcolumn 
\vspace{0.1cm} 
% Usamos el nuevo comando \CVTextBlock para el perfil
\CVTextBlock{Desarrollador de software y sistemas, enfocado en soluciones tecnológicas para la gestión administrativa. Comprometido con entregar resultados medibles mediante innovación accesible y trabajo en equipo multidisciplinario. Enfoque flexible y adaptativo, brindando productos escalables y sostenibles que responden a las necesidades específicas de cada proyecto.}

\end{paracol}

% =========================================================
% EXPERIENCIA LABORAL - SECCIÓN A ANCHO COMPLETO
% Título alineado a la izquierda (fuera de la columna).
% Columna izquierda: cargo / empresa / lugar / periodo.
% Columna derecha: descripción y logros.
% =========================================================
\CVSectionNormal{Experiencia}

\setlength{\columnsep}{1.6cm}
\begin{paracol}{2}
\columnratio{0.22}

% --- Entrada 1 ---
\CVExpMeta{Desarrollador Full-Stack}
{Independiente}
{Trabajo Remoto}
{OCT 2024 - ENE 2026}
\switchcolumn
\CVExpBody{
Soluciones digitales, aplicaciones móviles y sistemas comerciales digitales.
\begin{itemize}[label=\textcolor{CVBullet}{–}, leftmargin=0.5cm, topsep=0.4cm, itemsep=0.15cm]
    \item Implementación de métodos de pago seguros y sistema de apelaciones en plataforma de ventas digitales.
    \item Desarrollo de aplicación de redes sociales móviles con llamadas de voz y video integradas.
\end{itemize}
}
\switchcolumn*

% --- Entrada 2 ---
\CVExpMeta{Analista}
{DCTI / Gobernación del Estado Monagas}
{Maturín, Edo. Monagas}
{SEP 2023 - OCT 2024}
\switchcolumn
\CVExpBody{
Entidad pública dedicada a impulsar la modernización y eficiencia administrativa.
\begin{itemize}[label=\textcolor{CVBullet}{–}, leftmargin=0.5cm, topsep=0.4cm, itemsep=0.15cm]
    \item Creación de sistema integral para la organización y control de eventos institucionales.
    \item Desarrollo de API para la propuesta de un nuevo sistema de Plan Operativo Anual (POA).
    \item Diseño de módulos de usuario con accesos seguros y personalizados.
\end{itemize}
}
\switchcolumn*

% --- Entrada 3 ---
\CVExpMeta{Desarrollador Full-Stack}
{Independiente}
{Trabajo Remoto}
{ENE 2022 - JUL 2023}
\switchcolumn
\CVExpBody{
Desarrollo de soluciones tecnológicas personalizadas, enfocado en servicios integrales de desarrollo web y móvil.
\begin{itemize}[label=\textcolor{CVBullet}{–}, leftmargin=0.5cm, topsep=0.4cm, itemsep=0.15cm]
    \item Diseño e implementación de páginas web atractivas y funcionales con HTML y CSS.
    \item Creación de scripts inteligentes para automatizar tareas repetitivas en el servidor.
    \item Implementación de sistema seguro de inicio de sesión y acceso autorizado (OAuth2).
    \item Creación de API robusta y escalable para manejo de pedidos de tiendas en vivo.
\end{itemize}
}

\end{paracol}

% =========================================================
% CONTINUACIÓN DEL LAYOUT DIVIDIDO (Educación y Referencias)
% =========================================================
\newpage
% ENCABEZADO repetido (página 2)
\CVHeader
\vspace{1cm}
\begin{paracol}{2}

% --- EDUCACIÓN ---
\CVSection{Educación}
\switchcolumn
\vspace{0.1cm}

% Línea gris continua que abarca toda la sección
\noindent
\textcolor{CVLine}{\vrule width 1pt} \hspace{1cm}
\begin{minipage}[t]{\dimexpr\linewidth-0.4cm-1.5pt}
    \CVEdItem{Certificado de Iniciación al Desarrollo con IA}
    {Big School}
    {En línea}
    {MAR 2026}

    \CVEdItem{Técnico Superior Universitario en Informática}
    {Universidad Bolivariana de Venezuela}
    {Maturín, Edo. Monagas}
    {DIC 2025}

    \CVEdItem{Certificación en Algoritmos y Programación}
    {Universidad Bolivariana de Venezuela}
    {Maturín, Edo. Monagas}
    {FEB 2023}

    \CVEdItem{Certificación en Formación de Microempresas}
    {Programa "Simón Rodríguez" Núcleo Monagas}
    {Maturín, Edo. Monagas}
    {DIC 2021}
\end{minipage}

\end{paracol}

% =========================================================
% REFERENCIAS - SECCIÓN A ANCHO COMPLETO (título a la izquierda)
% =========================================================
\CVSectionNormal{Referencias}

\CVRef{Dr.C. Victor Rivas}{Doctorado en Ciencias Pedagógicas}{UBV}{+58 416 1819 143}
\hfill
\CVRef{Oduber Vasquez}{Profesional Universitario}{DCTI}{+58 424 9226 010}

\vspace{0.2cm}

% =========================================================
% SECCIONES MODO "NORMAL" (Ancho Completo)
% =========================================================

% --- SOFTWARE / HABILIDADES ---
\CVSectionNormal{Habilidades}
{\DMSansLight
\begin{multicols}{2}
    \textcolor{CVBullet}{–} Node.js \\
    \textcolor{CVBullet}{–} TypeScript \\
    \textcolor{CVBullet}{–} Dart / Flutter \\
    \textcolor{CVBullet}{–} Golang \\
    \textcolor{CVBullet}{–} PHP \\
    \textcolor{CVBullet}{–} SQL \\
    \textcolor{CVBullet}{–} HTML/CSS \\
    \textcolor{CVBullet}{–} JavaScript \\
    \textcolor{CVBullet}{–} UI/UX
\end{multicols}
}

% --- INTERESES EXTRACURRICULARES ---
\CVSectionNormal{Intereses Extracurriculares}
{\DMSansLight
\textcolor{CVBullet}{–} Idiomas \\
\textcolor{CVBullet}{–} Publicidad y mercadeo \\
\textcolor{CVBullet}{–} E-commerce \\
\textcolor{CVBullet}{–} Diseño \\
\textcolor{CVBullet}{–} Automatización de sistemas
}

\end{document}